\documentclass[10pt,a4paper]{amsart}
\usepackage[margin=19mm]{geometry}
\usepackage{amsmath,amssymb,mathtools}
\usepackage[scr=rsfs,cal=euler]{mathalfa}
\usepackage{xfrac}
\usepackage[unicode,hidelinks,bookmarks=false]{hyperref}
\newcommand{\ie}{i.e.\ }
\newcommand{\cf}{cf.\ }
\newcommand{\resp}{respectively\ }
\newcommand{\Subsection}{\subsection}
\providecommand{\nobreakdash}{\nobreak\hbox{-}}
\setlength{\emergencystretch}{2em}
\multlinegap0pt
\newtheorem{remark}{Remark}
\newtheorem{lemma}{Lemma}
\usepackage{cleveref}
\crefname{remark}{Remark}{Remarks}
\crefname{lemma}{Lemma}{Lemmas}
\title{On the expansion of Hanoi graphs}
\author{David Eppstein, Daniel Frishberg, William Maxwell}
\date{Source: arXiv:2510.18010v2 (CC BY 4.0)}
\begin{document}
\maketitle

\noindent\emph{Source and licence.} On the expansion of Hanoi graphs. Version arXiv:2510.18010v2 (CC BY 4.0), distributed by arXiv under the Creative Commons Attribution 4.0 International licence, http://creativecommons.org/licenses/by/4.0/, as recorded in the arXiv licence metadata for this exact version and shown on the abstract page https://arxiv.org/abs/2510.18010v2. Full licence notice: \texttt{licenses/arxiv-2510.18010-CC-BY-4.0.txt}.\par\smallskip

\noindent\textit{Editorial scope.} Counted unit: the article's lemma lem:hanpartialcong (source0.tex lines 575--578). The article's complete original proof of it is delivered with it (source0.tex lines 579--596, one \texttt{proof} environment of 18 lines). No mathematics is written, repaired or re-derived: the statement and the proof are the article's own lines, in the article's own notation, and no retained line is substituted or re-flowed.  Retained uncounted as the source-local context the proof needs: lines 532--534; lines 553--555.  Nothing was omitted from any retained span, so the package carries no omission record.  No retained line contains a citation, so the wrapper carries no bibliography and no bibliography entry.  No retained line contains a figure environment, an image-inclusion command or drawing code, so no image is delivered and the package declares no asset dependency.  Editorial formatting only, in addition: an AMS \texttt{amsart} preamble that restates the article's own declarations and macros used by the retained lines, namely the environments \texttt{remark}, \texttt{lemma} on separate counters, together with the standard packages those lines need.  The retained environments print in this document as Remark 1, Lemma 1 in order of appearance, whereas the article prints the same items with the numbering of its own full document.  The complete native source of the article as distributed in the arXiv e-print for this exact version is bundled unchanged as \texttt{sources/187465/source0.tex}.
\begin{remark}\label{rmk:fdistsubrecj}
Each of~$\pi_{distsub,1}$ and~$\pi_{distsub,2}$ is the sum of~$p-2$ recursive subproblems of the same form as~$\pi_{dist}$.
\end{remark}

\begin{lemma}\label{lem:fdistroutrec}
The subproblem~$\pi_{rout,i,j}$ can be solved with congestion~$\frac{\sigma_{rout,i,j}}{|V(H_p^n)|}$.
\end{lemma}

\begin{lemma}
\label{lem:hanpartialcong}
There exists an MSF~$f_{dist}$ that solves the MSF problem~$\pi_{dist}$ while producing at most~$\frac{\sigma_{dist}}{|V(H_p^n)|}$ congestion across each edge in~$H_2$.
\end{lemma}

\begin{proof}

The MSF~$f_{dist}$, given the decomposition, is entirely specified by the definitions we have given for the flows solving each subproblem.

We now apply induction, on $n$, to the recursive definition of~$f_{dist}$, to bound the congestion. The inductive hypothesis is the assumption that each~$\pi_{distsub,i}$ can be solved while producing at most~$\sigma_{distsub,i}$ (non-normalized) flow across each edge.

The distribution subproblems~$\pi_{distsub,i}$ then produce $\sigma_{distsub,i} = ((p-2)/p)\sigma_{dist}$ combined flow across edges within each~$H_{2,i}$,  $3 \leq i \leq p$\textemdash this is immediate from the inductive hypothesis.

Each of the subproblems $\pi_{distsub,1}$ and $\pi_{distsub,2}$ naturally decomposes (\cref{rmk:fdistsubrecj}) into $p-2$ subproblems, one for each of the boundary sets $\delta_v(H_{2,1}, H_{2,i})$, $i = 3, \dots, p$ and similarly one for each of the boundary sets $\delta_v(H_{2,2}, H_{2,i})$, $i = 3, \dots, p$. This produces total flow $\frac{(p-2)\sigma_{dist}}{p}$ across edges in~$H_{2,1}$ and~$H_{2,2}$ respectively.

Now, to bound the congestion resulting from~$\pi_{rout,i,j}$, recall that $\sigma_{rout,i,j} = \delta_{rout,i,j} = \sigma_{tran,i,j} = (1/p)\sigma_{dist}$. Applying \cref{lem:fdistroutrec} and summing over $j \in \{1, 2\}$, gives total (non-normalized) flow
$$(2/p)\sigma_{dist}$$
from the routing MSFs within each $H_{2,i}$. 

Adding the~$\pi_{rout,i,j}$ bound of~$(2/p)\sigma_{dist}$ to the $f_{distsub,i}$ congestion bound of $((p-2)/p)\sigma_{dist}$ shows that~$f_{dist}$ produces total congestion at most $\sigma_{dist}$ within each~$H_{2,i}$, $3 \leq i \leq p$, as claimed.

The claim now follows from normalizing the above congestion bounds by $|V(H_p^n)|$.
\end{proof}

\end{document}
